\documentclass{ximera}

\ifdefined\HCode
  \newcommand{\alttext}[1]{\HCode{<span class="sr-only">#1</span>}}
\else
  \newcommand{\alttext}[1]{} % Fallback for PDF view
\fi


% MATH -----------------------------------------------------------
\newcommand{\nats}{\mathbb N}
\newcommand{\ints}{\mathbb Z}
\newcommand{\rats}{\mathbb Q}
\newcommand{\reals}{\mathbb R}
\newcommand{\complex}{\mathbb C}
\newcommand{\powerset}{\mathscr P}

%-------------------------------------------------------------

\pgfplotsset{compat=1.18}
\begin{document}
\begin{abstract}
 \end{abstract}

    \title{Exemplar Examples 23}
    
    \maketitle

A $10$-kg mass is hung on a very large undamped vertical spring and brought to rest at equilibrium, stretching the spring by $39.2$ centimeters.  Then an external upward force of $F(t)=10t$ N is applied at time $0\leq t<10$ seconds, followed by no external force (free motion).


\begin{enumerate}
\item Using $g=9.8$ m/s$^2$ for the magnitude of the acceleration due to gravity, let's determine the spring stiffness $k$ in N/m.

% k=mg/l=98/0.392=250 N/m

\vfill

\item Let's derive an IVP satisfied by $y(t)$, the displacement of the mass above equilibrium (in meters) at time $t\geq 0$ seconds.

\vfill

%   $10 y'' + 250y= 10t - U(t-10) [10t] , y(0)=0, y'(0)=0

\newpage

 \item Let's solve this IVP and give the solution as a piecewise defined function.


% Re-write as y''+25y=t -  U(t-10) [t]

% (s^2+25) Y = 1/s^2 - e^(-10s) [1/s^2 +10/s]

% Y = 1/[s^2(s^2+25)] - e^(-10s) [1/[s^2(s^2+25) ] +10/[s(s^2+25)]]


% PFD:  1/[s^2(s^2+25)] = (1/25) /s^2  - [1/25]/[s^2+25] and 10/[s(s^2+25)]=(2/5)/s - (2/5)s/(s^2+25)


% Y = (1/25) /s^2  - [1/25]/[s^2+25] - e^(-10s) [(1/25) /s^2  - [1/25]/[s^2+25] + (2/5)/s - (2/5)s/(s^2+25) ]

% y= t/25 -1/125 sin(5t) - U(t-10) [ (t-10)/25 -1/125 sin(5(t-10)) + 2/5 -2/5 cos(5(t-10))]

% y= t/25 -1/125 sin(5t) - U(t-10) [ t/25 -1/125 sin(5(t-10)) -2/5 cos(5(t-10))]

\vfill

\begin{figure}[h]
\alttext{The displacement of a mass from equilibrium in a Spring-mass system with piecewise continuous forcing.}
\centering
\begin{tikzpicture}
\begin{axis}[
    axis lines = middle,
    xlabel = {$t$},
    ylabel = {$y$},
    domain = 0:20,
    samples = 200,
    grid = both,
    grid style = {dashed, gray!30},
    xmin = 0, xmax = 20,
    ymin = -0.6, ymax = 1.2,
    title = {Graph of Piecewise Function $y(t)$},
    legend pos = north west
]

% Define the first part of the function (for t < 10)
\addplot [
    blue, 
    thick, 
    domain=0:10
] {x/25 - (1/125)*sin(deg(5*x))};

% Define the second part of the function (for t >= 10)
% We subtract the U(t-10) terms from the original function
\addplot [
    black, 
    thick, 
    domain=10:20
] {
    (x/25 - (1/125)*sin(deg(5*x))) - 
    (x/25 - (1/125)*sin(deg(5*(x-10))) - (2/5)*cos(deg(5*(x-10))))
};

\legend{$t < 10$, $t \ge 10$}

\end{axis}
\end{tikzpicture}
\caption{Spring-mass system with piecewise continuous forcing}
\end{figure}

\end{enumerate}
    
\end{document}